\documentclass{beamer}

\usetheme{Madrid}

\usepackage[spanish,es-noshorthands]{babel}
\usepackage[utf8]{inputenc}
\usepackage[T1]{fontenc}
\usepackage{amsmath,amssymb}

%------------------------------------------------
\title{Pensamiento Matemático}
\subtitle{Unidad II: Operatoria Básica}
\author{Profesor Luis Tulio González Alcaino \\ Magíster en Matemática}
\institute{Universidad Santo Tomás \\ Departamento Ciencias Básicas - Talca}
\date{Semestre I - 2026}

%------------------------------------------------
\begin{document}
	
	%------------------------------------------------
	\begin{frame}
		\titlepage
	\end{frame}
	
	%------------------------------------------------
	\begin{frame}{Contenidos de la clase}
		\begin{enumerate}
			\item Conjuntos Numéricos
			\item Operatoria
			\item Patrones e inconsistencias
		\end{enumerate}
	\end{frame}
	
	%------------------------------------------------
	\begin{frame}{Resultado de Aprendizaje}
		\begin{enumerate}
			\item Aplica operatoria con números enteros y racionales.
			\item Integra patrones y relaciones matemáticas.
			\item Detecta inconsistencias en resultados matemáticos.
		\end{enumerate}
	\end{frame}
	
	%------------------------------------------------
	\begin{frame}{Conjuntos Numéricos}
		\begin{enumerate}
			\item \textbf{Números Naturales \(\mathbb{N}\)}
			\[
			\mathbb{N}=\{1,2,3,\dots\}
			\]
			
			\pause
			
			\item \textbf{Números Enteros \(\mathbb{Z}\)}
			\[
			\mathbb{Z}=\{\dots,-2,-1,0,1,2,\dots\}
			\]
		\end{enumerate}
	\end{frame}
	
	%------------------------------------------------
	\begin{frame}{Conjuntos Numéricos}
		\begin{enumerate}
			\setcounter{enumi}{2}
			
			\item \textbf{Números Racionales \(\mathbb{Q}\)}
			\[
			\mathbb{Q}=\left\{\frac{a}{b}: b\neq0\right\}
			\]
			
			Ejemplos:
			\[
			\frac{1}{4}=0.25,\quad \frac{5}{3}=1.666\ldots
			\]
			
			\pause
			
			\item \textbf{Números Irracionales}
			\[
			\sqrt{2},\ \pi,\ e
			\]
		\end{enumerate}
	\end{frame}
	
	%------------------------------------------------
	\begin{frame}{Números Reales}
		\[
		\mathbb{R}=\mathbb{Q}\cup \mathbb{I}
		\]
		
		Representación en la recta numérica:
		
		\[
		\text{---}\quad -2 \quad -1 \quad 0 \quad 1 \quad 2 \quad \text{---}
		\]
	\end{frame}
	
	%------------------------------------------------
	\begin{frame}{Adición de enteros}
		\begin{itemize}
			\item Mismo signo: se suman y se conserva el signo
			\item Distinto signo: se restan y se conserva el signo mayor
		\end{itemize}
		
		\pause
		
		Ejemplos:
		\[
		(-4)+(-7)=-11
		\]
		
		\[
		8+(-3)=5
		\]
	\end{frame}
	
	%------------------------------------------------
	\begin{frame}{Observaciones}
		\begin{enumerate}
			\item \(a+(-b)=a-b\)
			\item \(a-(-b)=a+b\)
			\item \(a+0=a\)
			\item \(a+(-a)=0\)
			\item \(-(a+b)=(-a)+(-b)\)
		\end{enumerate}
	\end{frame}
	
	%------------------------------------------------
	\begin{frame}{Ejercicios}
		\begin{enumerate}
			\item \((-3)+(-8)\)
			\item \((-5+3)+2\)
			\item \(10+(-5)\)
			\item \(0+(-5+2)\)
			\item \((-32)+48+(-10)\)
		\end{enumerate}
	\end{frame}
	
	%------------------------------------------------
	\begin{frame}{Multiplicación}
		Ley de signos:
		
		\[
		(+)(+)=+ \quad (-)(-) = +
		\]
		
		\[
		(+)(-)=- \quad (-)(+)=- 
		\]
		
		Ejemplos:
		\[
		(-3)(-7)=21
		\]
		\[
		4(-5)=-20
		\]
	\end{frame}
	
	%------------------------------------------------
	\begin{frame}{Jerarquía de operaciones}
		\begin{enumerate}
			\item Paréntesis
			\item Potencias
			\item Multiplicación y división
			\item Suma y resta
		\end{enumerate}
	\end{frame}
	
	%------------------------------------------------
	\begin{frame}{Ejemplo}
		\[
		-3+2(5-8)
		\]
		
		\pause
		
		\[
		-3+2(-3)
		\]
		
		\pause
		
		\[
		-3-6=-9
		\]
	\end{frame}
	
	%------------------------------------------------
	\begin{frame}{Números racionales}
		Suma:
		\[
		\frac{a}{b}+\frac{c}{b}=\frac{a+c}{b}
		\]
		
		\pause
		
		Ejemplo:
		\[
		\frac{2}{3}+\frac{1}{6}=\frac{5}{6}
		\]
	\end{frame}
	
	%------------------------------------------------
	\begin{frame}{Multiplicación de racionales}
		\[
		\frac{a}{b}\cdot\frac{c}{d}=\frac{ac}{bd}
		\]
		
		Ejemplo:
		\[
		\frac{3}{5}\cdot\frac{10}{9}=\frac{2}{3}
		\]
	\end{frame}
	
	%------------------------------------------------
	\begin{frame}{División de racionales}
		\[
		\frac{a}{b}\div\frac{c}{d}=\frac{a}{b}\cdot\frac{d}{c}
		\]
		
		Ejemplo:
		\[
		\frac{4}{7}\div\frac{2}{3}=\frac{6}{7}
		\]
	\end{frame}
	
	%------------------------------------------------
	\begin{frame}{Patrones}
		\[
		2,\ 5,\ 8,\ 11,\dots
		\]
		
		Patrón: sumar 3
		
		\pause
		
		\[
		\frac12,\ 1,\ \frac32,\ 2,\dots
		\]
		
		Patrón: sumar \(\frac12\)
	\end{frame}
	
	%------------------------------------------------
	\begin{frame}{Inconsistencias}
		Ejemplos incorrectos:
		
		\[
		-6+2=-8
		\]
		
		\pause
		
		\[
		\frac{1}{2}+\frac{1}{3}=\frac{2}{5}
		\]
		
		\pause
		
		Correcciones:
		\[
		-6+2=-4
		\]
		
		\[
		\frac{1}{2}+\frac{1}{3}=\frac{5}{6}
		\]
	\end{frame}
	
	%------------------------------------------------
	\begin{frame}{Aplicaciones}
		\begin{enumerate}
			\item Variación:
			\[
			+5-3+4-2+6
			\]
			
			\item Fracciones:
			\[
			\frac{3}{4}-\frac{1}{8}
			\]
			
			\item Patrón:
			\[
			10,\ 8,\ 6,\ 4,\dots
			\]
		\end{enumerate}
	\end{frame}
	
	%------------------------------------------------
	\begin{frame}{Cierre}
		\begin{itemize}
			\item Operatoria con enteros
			\item Operatoria con racionales
			\item Patrones
			\item Detección de errores
		\end{itemize}
		
		\vspace{0.3cm}
		
		\centering
		\textbf{No basta con calcular, hay que analizar y verificar.}
	\end{frame}
	
	%------------------------------------------------
	\begin{frame}
		\centering
		\Huge Gracias
	\end{frame}
	
\end{document}